\documentclass[aspectratio=169,11pt,usenames,dvipsnames,handout]{beamer}

\usepackage[czech,english]{babel}
\usepackage{CJKutf8}
\newcommand{\zh}[1]{\begin{CJK}{UTF8}{gbsn}#1\end{CJK}}
\usepackage{graphicx}
\usepackage{enumitem}
\usepackage{amsmath}
\newlength{\dhatheight}
\newcommand{\doublehat}[1]{%
    \settoheight{\dhatheight}{\ensuremath{\hat{#1}}}%
    \addtolength{\dhatheight}{-0.2ex}%
    \hat{\vphantom{\rule{1pt}{\dhatheight}}%
    \smash{\hat{#1}}}}
\usepackage{mathtools}
\usepackage{float}
\usepackage{tikz}
\usetikzlibrary{patterns,arrows.meta,calc,3d,angles}
\usepackage{tkz-euclide}
\tikzset{point style/.style = {%
  draw = black,
  inner sep = 0pt,
  shape = circle,
  minimum size = 5pt,
  fill = black
 }
}
\usepackage{enumitem}

\usepackage{caption}
\usepackage{subcaption}

\usepackage{booktabs}
% Flowchart stuff

\usepackage{pgfopts}
\usepackage{xcolor}
\usepackage{tcolorbox}

\usetheme[
 titlestyle=style2,
 titleformat=smallcaps,
 sectionstyle=plain,
 slidestyle=cyber,
 headingcolor=theme,
 block=transparent
]{trigon}

\title{Systems of Linear Equations}
\date{\today}
\author{Adam Klepáč}
\institute[GEVO]{Gymnázium Evolution Jižní Město}
\biglogo[width=.2\textwidth]{logo}
\smalllogo[width=.1\textwidth]{logo}
\titlegraphic{\includegraphics[height=\paperheight]{title}}

\def\subsectionname{}

% enumerate and itemize global settings
\setlist{topsep=0pt}
\setlist[itemize,1]{label=\textbullet}
\setlist[itemize,2]{label=\textopenbullet}
\setlist[enumerate,1]{label=\arabic*.}
\setlist[enumerate,2]{label=\alph*)}

% custom colors %
\definecolor{RichBlack}{HTML}{050F1F}
\definecolor{PowderBlue}{HTML}{9ABBD8}
\definecolor{BlueGray}{HTML}{6897BE}
\definecolor{BerkeleyBlue}{HTML}{092B55}
\definecolor{LapisLazuli}{HTML}{40698F}

\colorlet{tPrim}{BerkeleyBlue}
\colorlet{tTheme}{PowderBlue}
\colorlet{tSec}{LapisLazuli}
\colorlet{tAccent}{BlueGray}
\colorlet{tText}{RichBlack}

\newcommand{\clr}{\textcolor{BrickRed}}
\newcommand{\clb}{\textcolor{RoyalBlue}}
\newcommand{\clg}{\textcolor{ForestGreen}}
\newcommand{\clm}{\textcolor{Magenta}}
\newcommand{\cls}{\textcolor{Salmon}}
\newcommand{\clo}{\textcolor{OrangeRed}}

\newcommand{\N}{\mathbb{N}}
\newcommand{\Z}{\mathbb{Z}}
\newcommand{\R}{\mathbb{R}}
\renewcommand{\P}{\mathbb{P}}
\DeclareMathOperator{\tng}{\triangle}
\DeclareMathOperator{\dv}{div}

\tcbset{
 boxsep=7pt,
 fonttitle=\sc,
 colframe=tGreyBg,
 colframe=tSec,
 boxrule=1pt
}

\begin{document}
\titleframe

\begin{frame}
 \frametitle{Contents}
 \tableofcontents
\end{frame}

\section{Functions}

\begin{frame}
 \frametitle{What Is A Function?}
 Intuitively, a function is a \alert{box} which receives data and gives some
 data back.\pause
 \begin{center}
  \begin{tikzpicture}
   \node[draw,rectangle,minimum height=1cm,minimum
    width=2cm,color=BlueGray,align=center] (f)
    at (0,0) {function \\ \footnotesize `\emph{do something with \clr{input}'}};
   \node[left=2cm of f] (in) {\clr{input}};
   \node[right=2cm of f] (out) {\clb{output}};
   \draw[-latex,thick,shorten <= 2pt, shorten >= 2pt] (in) -- (f);
   \draw[-latex,thick,shorten <= 2pt, shorten >= 2pt] (f) -- (out);
  \end{tikzpicture}
 \end{center}
 \pause
 We'll call the data that \alert{a function receives}, \clr{inputs} and the
 \alert{data it gives back}, \clb{outputs}. \\ \pause
 \clr{Inputs} and \clb{outputs} need not necessarily be just `one object', they
 can be for example lists of numbers.
\end{frame}

\begin{frame}
 \frametitle{Functions -- Example}
 A function which returns the \alert{average} of a given set of numbers receives
 the numbers and also their count as \clr{input} and returns the \alert{average}
 as \clb{output}.\pause
 \begin{center}
  \begin{tikzpicture}
   \node[draw,rectangle,minimum height=1cm,minimum
    width=2cm,color=BlueGray,align=center] (f)
    at (0,0) {average \\ \footnotesize `\emph{sum all the numbers}\\
     \footnotesize \emph{and
    divide them by $n$}'};
   \node[left=2cm of f,align=center] (in)
    {\clr{$x_1,x_2,\ldots,x_n$}\\\clr{$n$}}; \node[right=2cm of f] (out)
     {\clb{$\displaystyle \frac{x_1+x_2+\ldots +x_n}{n}$}};
   \draw[-latex,thick,shorten <= 2pt, shorten >= 2pt] (in) -- (f);
   \draw[-latex,thick,shorten <= 2pt, shorten >= 2pt] (f) -- (out);
  \end{tikzpicture}
 \end{center}
\end{frame}

\begin{frame}
 \frametitle{Functions -- Example}
 We can also consider `non-mathematical' functions. Like a function which
 receives a type of meal and returns the ingredients.\\ \pause
 \begin{center}
  \begin{tikzpicture}
   \node[draw,rectangle,minimum height=1cm,minimum
    width=2cm,color=BlueGray,align=center] (f) at (0,0) {ingredients};
   \node[left=2cm of f,align=center] (in) {\clr{omelette}};
   \node[right=2cm of f,align=center] (out) {\clb{eggs} \\ \clb{salt} \\
    \clb{butter}};
   \draw[-latex,thick,shorten <= 2pt, shorten >= 2pt] (in) -- (f);
   \draw[-latex,thick,shorten <= 2pt, shorten >= 2pt] (f) -- (out);
  \end{tikzpicture}
 \end{center}
\end{frame}

\subsection{Function Composition}

\begin{frame}
 \subsectionpage
\end{frame}

\begin{frame}
 \frametitle{Function Composition}
 If we have \alert{two functions}, we can \alert{in certain cases} `compose'
 them. \\ \pause
 \alert{Composition} simply means that \alert{one function follows the other} --
 in other words, the \clb{output} of the first function is the \clr{input} of
 the second.\\ \pause
 Of course, \alert{composition} is only possible if the \clb{output} of the
 first function is a valid \clr{input} for the second. \\ \pause
 For instance, you could hardly compose the \alert{ingredients} function with
 the \alert{average} function.
\end{frame}

\begin{frame}
 \frametitle{Function Composition}
 Considering two functions
 \vspace*{-1em}
 \begin{center}
  \begin{tikzpicture}
   \node[draw,rectangle,minimum height=.5cm,minimum
    width=2cm,color=BlueGray,align=center] (f) at (0,0) {sum};
   \node[left=2cm of f,align=center] (in) {\clr{$a,b$}};
   \node[right=2cm of f,align=center] (out) {\clb{$a + b$}};
   \draw[-latex,thick,shorten <= 2pt, shorten >= 2pt] (in) -- (f);
   \draw[-latex,thick,shorten <= 2pt, shorten >= 2pt] (f) -- (out);

   \node[draw,rectangle,minimum height=.5cm,minimum
    width=2cm,color=BlueGray,align=center,below=.5cm of sum] (g) {times two};
   \node[left=2cm of g,align=center] (in) {\clr{$a$}};
   \node[right=2cm of g,align=center] (out) {\clb{$2 \cdot a$}};
   \draw[-latex,thick,shorten <= 2pt, shorten >= 2pt] (in) -- (g);
   \draw[-latex,thick,shorten <= 2pt, shorten >= 2pt] (g) -- (out);
  \end{tikzpicture}
 \end{center}
 \pause
 their composition can look like this
 \vspace*{-1em}
 \begin{center}
  \begin{tikzpicture}
   \node[draw,rectangle,minimum height=.5cm,minimum
    width=2cm,color=BlueGray,align=center] (f) at (0,0) {sum};
   \node[draw,rectangle,minimum height=.5cm,minimum
    width=2cm,color=BlueGray,align=center,right=1cm of f] (g) {times two};
   \node[left=1cm of f,align=center] (in) {\clr{$a,b$}};
   \node[right=1cm of g,align=center] (out) {\clb{$2 \cdot (a + b)$}};
   \draw[-latex,thick,shorten <= 2pt, shorten >= 2pt] (in) -- (f);
   \draw[-latex,thick,shorten <= 2pt, shorten >= 2pt] (f) -- (g);
   \draw[-latex,thick,shorten <= 2pt, shorten >= 2pt] (g) -- (out);
  \end{tikzpicture}
 \end{center}
 \pause
 What would the output of this composition look like
 \begin{center}
  \begin{tikzpicture}
   \node[draw,rectangle,minimum height=.5cm,minimum
    width=2cm,color=BlueGray,align=center] (f) at (0,0) {times two};
   \node[draw,rectangle,minimum height=.5cm,minimum
    width=2cm,color=BlueGray,align=center,right=1cm of f] (g) {sum};
   \node[left=1cm of f,align=center] (in) {\clr{$a$}};
   \node[right=1cm of g,align=center] (out) {\clb{?}};
   \draw[-latex,thick,shorten <= 2pt, shorten >= 2pt] (in) -- (f);
   \draw[-latex,thick,shorten <= 2pt, shorten >= 2pt] (f) -- (g);
   \draw[-latex,thick,shorten <= 2pt, shorten >= 2pt] (g) -- (out);
  \end{tikzpicture}
 \end{center}
\end{frame}

\begin{frame}
 \frametitle{Function Composition}
 So, the \alert{order of the composition matters}! Here are a few examples of
 compositions which make sense:
 \begin{center}
  \begin{tikzpicture}
   \node[draw,rectangle,minimum height=.5cm,minimum
    width=2cm,color=BlueGray,align=center] (f) at (0,0) {ingredients};
   \node[draw,rectangle,minimum height=.5cm,minimum
    width=2cm,color=BlueGray,align=center,right=1cm of f] (g) {total price};
   \node[left=1cm of f,align=center] (in) {\clr{omelette}};
   \node[right=1cm of g,align=center] (out) {\clb{\$4.37}};
   \draw[-latex,thick,shorten <= 2pt, shorten >= 2pt] (in) -- (f);
   \draw[-latex,thick,shorten <= 2pt, shorten >= 2pt] (f) -- (g);
   \draw[-latex,thick,shorten <= 2pt, shorten >= 2pt] (g) -- (out);
  \end{tikzpicture}\\
  \pause
  \vspace*{1em}
  \begin{tikzpicture}
   \node[draw,rectangle,minimum height=.5cm,minimum
    width=2cm,color=BlueGray,align=center] (f) at (0,0) {average};
   \node[draw,rectangle,minimum height=.5cm,minimum
    width=2cm,color=BlueGray,align=center,right=1cm of f] (g) {times two};
   \node[left=1cm of f,align=center] (in) {\clr{$x_1,x_2,\ldots,x_n$}\\\clr{$n$}};
   \node[right=1cm of g,align=center] (out) {\clb{$\displaystyle 2 \cdot
    \frac{x_1+x_2+\ldots+x_n}{n}$}};
   \draw[-latex,thick,shorten <= 2pt, shorten >= 2pt] (in) -- (f);
   \draw[-latex,thick,shorten <= 2pt, shorten >= 2pt] (f) -- (g);
   \draw[-latex,thick,shorten <= 2pt, shorten >= 2pt] (g) -- (out);
  \end{tikzpicture}
 \end{center}
 \pause
 We can of course compose \alert{as many functions} as we like. An example of
 this:
 \begin{center}
  \begin{tikzpicture}
   \node[draw,rectangle,minimum height=.5cm,minimum
    width=2cm,color=BlueGray,align=center] (f) at (0,0) {ingredients};
   \node[draw,rectangle,minimum height=.5cm,minimum
    width=2cm,color=BlueGray,align=center,right=1cm of f] (g) {total price};
   \node[draw,rectangle,minimum height=.5cm,minimum
    width=2cm,color=BlueGray,align=center,right=1cm of g] (h) {times two};
   \node[left=1cm of f,align=center] (in) {\clr{omelette}};
   \node[right=1cm of h,align=center] (out) {\clb{\$8.74}};
   \draw[-latex,thick,shorten <= 2pt, shorten >= 2pt] (in) -- (f);
   \draw[-latex,thick,shorten <= 2pt, shorten >= 2pt] (f) -- (g);
   \draw[-latex,thick,shorten <= 2pt, shorten >= 2pt] (g) -- (h);
   \draw[-latex,thick,shorten <= 2pt, shorten >= 2pt] (h) -- (out);
  \end{tikzpicture}\\
 \end{center}
\end{frame}

\begin{frame}
 \frametitle{Functions -- Notation}
 Drawing pictures like this would be cumbersome. \pause Instead of
 \begin{center}
  \begin{tikzpicture}
   \node[draw,rectangle,minimum height=1cm,minimum
    width=2cm,color=BlueGray,align=center] (f) at (0,0) {function};
   \node[left=2cm of f,align=center] (in) {\clr{input}};
   \node[right=2cm of f,align=center] (out) {\clb{output}};
   \draw[-latex,thick,shorten <= 2pt, shorten >= 2pt] (in) -- (f);
   \draw[-latex,thick,shorten <= 2pt, shorten >= 2pt] (f) -- (out);
  \end{tikzpicture}
 \end{center}
 we simply write \alert{function}(\clr{input}) = \clb{output}. \\ \pause
 You're probably used to seeing function written like $f(x) = y$. The picture
 corresponding to this is
 \begin{center}
  \begin{tikzpicture}
   \node[draw,rectangle,minimum height=1cm,minimum
    width=1cm,color=BlueGray,align=center] (f) at (0,0) {$f$};
   \node[left=2cm of f,align=center] (in) {\clr{$x$}};
   \node[right=2cm of f,align=center] (out) {\clb{$y$}};
   \draw[-latex,thick,shorten <= 2pt, shorten >= 2pt] (in) -- (f);
   \draw[-latex,thick,shorten <= 2pt, shorten >= 2pt] (f) -- (out);
  \end{tikzpicture}
 \end{center}
\end{frame}

\begin{frame}
 \frametitle{Functions -- Notation}
 The symbol used for function composition is $ \circ $. It is however a little
 confusing because it is written in order `from right to left'. \\ \pause
 For example, if $f$ and $g$ are two functions, their composition $f \circ g$
 corresponds to this picture
 \begin{center}
  \begin{tikzpicture}
   \node[draw,rectangle,minimum height=1cm,minimum
    width=1cm,color=BlueGray,align=center] (g) at (0,0) {$g$};
   \node[draw,rectangle,minimum height=1cm,minimum
    width=1cm,color=BlueGray,align=center,right=2cm of g] (f) at (0,0) {$f$};
   \node[left=1cm of g,align=center] (in) {\clr{$x$}};
   \node[right=1cm of f,align=center] (out) {\clb{$y$}};
   \draw[-latex,thick,shorten <= 2pt, shorten >= 2pt] (in) -- (g);
   \draw[-latex,thick,shorten <= 2pt, shorten >= 2pt] (g) -- (f);
   \draw[-latex,thick,shorten <= 2pt, shorten >= 2pt] (f) -- (out);
  \end{tikzpicture}
 \end{center}
 that is, \alert{first $g$, then $f$}.
\end{frame}

\subsection{Real Functions}

\begin{frame}
 \subsectionpage
\end{frame}

\begin{frame}
 \frametitle{Real Function}
 A \alert{real function} is simply a function whose \alert{\clr{input} and
 \clb{output} are both real numbers}. \\ \pause
 Examples of such functions are
 \begin{itemize}
  \item $f(x) = 0$,
  \item $g(x) = \tan^{6}(\log^{\sin(x^2 + 4)}(\frac{5x^3 - 2}{9x^{7}}))$,
 \end{itemize}
 where $x \in \R$. \pause Or, using pictures,
 \begin{center}
  \begin{tikzpicture}
   \node[draw,rectangle,minimum height=1cm,minimum
    width=1cm,color=BlueGray,align=center] (f) at (0,0) {$f$};
   \node[left=2cm of f,align=center] (in) {\clr{$x$}};
   \node[right=2cm of f,align=center] (out) {\clb{$0$}};
   \draw[-latex,thick,shorten <= 2pt, shorten >= 2pt] (in) -- (f);
   \draw[-latex,thick,shorten <= 2pt, shorten >= 2pt] (f) -- (out);

   \node[draw,rectangle,minimum height=1cm,minimum
    width=1cm,color=BlueGray,align=center,below=.5cm of f] (g) {$g$};
   \node[left=2cm of g,align=center] (in) {\clr{$x$}};
   \node[right=2cm of g,align=center] (out) {\clb{$\tan^{6}(\log^{\sin(x^2 +
    4)}(\frac{5x^3 - 2}{9x^{7}}))$}};
   \draw[-latex,thick,shorten <= 2pt, shorten >= 2pt] (in) -- (g);
   \draw[-latex,thick,shorten <= 2pt, shorten >= 2pt] (g) -- (out);
  \end{tikzpicture}
 \end{center}
\end{frame}

\begin{frame}
 \frametitle{Operations On Real Functions}
 As both the \clr{input} and the \clb{output} of a real function is a real
 number, \alert{we can always compose real functions}. \\ \pause
 However, that \alert{doesn't mean that the order doesn't matter!} Different
 order of composition gives different functions.\\ \pause
 For example, take $f(x) = 2x^2 + 7$ and $g(x) = \frac{1}{1 + x}$. Then,\pause
 \begin{itemize}
  \item $(f \circ g)(x) = 2 \left( \frac{1}{1+x} \right)^2 + 7$ and \pause
  \item $(g \circ f)(x) = \frac{1}{1 + (2x^2 + 7)}$.
 \end{itemize}
\end{frame}

\begin{frame}
 \frametitle{Operations On Real Functions}
 Real functions can also be \alert{added} and \alert{multiplied}, just like real
 numbers.\\ \pause
 This involves simply adding or multiplying their respective outputs. \pause For
 the functions $f(x) = 2x^2 + 7$ and $g(x) = \frac{1}{1+x}$, their
 \begin{itemize}
  \item \alert{sum} is the function with output
  \[
   (f+g)(x) = f(x) + g(x) = 2x^2 + 7 + \frac{1}{1+x}.
  \]
  \pause
  \item \alert{product} is the function with output
  \[
   (f \cdot g)(x) = f(x) \cdot g(x) = (2x^2 + 7) \cdot \left( \frac{1}{1+x}
   \right) = \frac{2x^2+7}{1+x}.
  \]
 \end{itemize}
\end{frame}

\begin{frame}
 \frametitle{Graphs}
 As real functions have real numbers as inputs and outputs, they can be easily
 \alert{graphed}.\\ \pause
 \alert{Graphing} a real function $f$ simply means drawing the points
 $(\clr{x},\clb{f(x)})$ or $(\clr{input},\clb{output})$ in some chosen
 coordinate system.\\ \pause
 We typically use the \alert{Cartesian} coordinate system with two axes (one for
 \clr{input} and one for \clb{output}) that are mutually perpendicular. \pause
 These are often called the $\clr{x}$-axis and the $\clb{y}$-axis.\\ \pause
 However, later, we'll also use the \alert{polar} coordinate system where every
 point is instead determined by its angle and distance from the origin of the
 system.
\end{frame}

\begin{frame}
 \frametitle{Graphs}
 The functions $\clm{f(x) = 2x^2 + 7}$ and $\clg{g(x) = \frac{1}{1+x}}$ have the
 following (parts of) graphs:
 \begin{figure}[H]
  \centering
  \begin{subfigure}[b]{.48\textwidth}
   \centering
   \begin{tikzpicture}[scale=0.3]
    \tkzInit[xmin=-3,xmax=5,ymin=0,ymax=12]
    \tkzDrawX[>=latex,color=BrickRed]
    \tkzDrawY[>=latex,color=RoyalBlue]
    \draw[scale=0.5,color=Magenta,domain=-3:6,variable=\x,thick] plot
     ({\x},{0.5*\x*\x + 7}) node[right] {$\clm{f}$};
    \tkzDefPoints{2/0/a,0/7.5/fa,2/7.5/x}
    \tkzDrawSegments[dashed](a,x fa,x)
    \tkzDrawPoint[color=BrickRed](a)
    \tkzLabelPoint[below=1mm,color=BrickRed](a){$a$}
    \tkzDrawPoint[color=RoyalBlue](fa)
    \tkzLabelPoint[left,color=RoyalBlue](fa){$f(a)$}
    \tkzDrawPoint[color=Magenta](x)
    \tkzLabelPoint[right,color=Magenta](x){$(a,f(a))$}
   \end{tikzpicture}
  \end{subfigure}
  \begin{subfigure}[b]{.48\textwidth}
   \centering
   \begin{tikzpicture}[scale=0.3]
    \tkzInit[xmin=-3,xmax=10,ymin=0,ymax=12]
    \tkzDrawX[>=latex,color=BrickRed]
    \tkzDrawY[>=latex,color=RoyalBlue]
    \draw[scale=2,smooth,color=ForestGreen,domain=-0.8:4,variable=\x,thick] plot
     ({\x},{1/(1+\x)}) node[above] {$\clg{g}$};
    \tkzDefPoints{-1/0/a,0/4/fa,-1/4/x}
    \tkzDrawSegments[dashed](a,x fa,x)
    \tkzDrawPoint[color=BrickRed](a)
    \tkzLabelPoint[below=1mm,color=BrickRed](a){$a$}
    \tkzDrawPoint[color=RoyalBlue](fa)
    \tkzLabelPoint[right,color=RoyalBlue](fa){$g(a)$}
    \tkzDrawPoint[color=ForestGreen](x)
    \tkzLabelPoint[left,color=ForestGreen](x){$(a,g(a))$}
   \end{tikzpicture}
  \end{subfigure}
 \end{figure}
\end{frame}

\end{document} 
